\documentclass[aspectratio=169,11pt]{beamer}
\usetheme{default}
\usecolortheme{default}
\usepackage[utf8]{inputenc}
\usepackage[spanish]{babel}
\usepackage{graphicx}
\usepackage{tikz}
\usetikzlibrary{shapes.geometric,arrows.meta,positioning}
\usepackage{fontspec}
\setmainfont{DejaVu Sans}

\definecolor{verdehoja}{RGB}{76, 140, 74}
\definecolor{marronlarva}{RGB}{139, 90, 43}
\definecolor{grisfondo}{RGB}{250, 248, 245}
\definecolor{azulcielo}{RGB}{100, 149, 237}

\setbeamercolor{background canvas}{bg=grisfondo}
\setbeamercolor{frametitle}{fg=white,bg=verdehoja}
\setbeamercolor{title}{fg=marronlarva}
\setbeamercolor{subtitle}{fg=verdehoja}
\setbeamercolor{author}{fg=marronlarva}
\setbeamercolor{item}{fg=verdehoja}
\setbeamercolor{block title}{fg=white,bg=marronlarva}
\setbeamercolor{block body}{bg=white,fg=black}
\setbeamertemplate{navigation symbols}{}

\title{\textbf{La Mosca Soldado Negra:}\\De los residuos a la proteina}
\subtitle{Modelado inteligente de emisiones para una bioconversion sostenible}
\author{Dr.~(Cand.) John Jairo Leal Gomez\\{\small Profesor Asociado, UNAL Palmira}}
\date{}

\begin{document}

\begin{frame}
  \vspace{1cm}
  \centering
  {\Huge\textcolor{verdehoja}{$\bullet$}}\\[0.3cm]
  {\LARGE\textbf{La Mosca Soldado Negra}}\\[0.2cm]
  {\large De los residuos a la proteina}\\[0.5cm]
  {\normalsize\textcolor{marronlarva}{Modelado inteligente de emisiones para una bioconversion sostenible}}\\[1cm]
  {\small Dr.~(Cand.) John Jairo Leal Gomez}\\[0.1cm]
  {\footnotesize Profesor Asociado, Universidad Nacional de Colombia -- Palmira}\\[0.3cm]
  {\footnotesize\textcolor{gray}{Sistema Inteligente para el Modelado y Prediccion de Emisiones de GEI}}
\end{frame}

\begin{frame}{Que es \textit{Hermetia illucens}?}
  \begin{columns}
    \begin{column}{0.55\textwidth}
      \begin{itemize}
        \item Una mosca neotropical, inofensiva, que no pica ni transmite enfermedades.
        \item Sus \textbf{larvas} son los verdaderos protagonistas: voraces, resistentes, eficientes.
        \item En 14 dias, una larva puede multiplicar su peso \textbf{100 veces}.
        \item Lo hacen \textbf{comiendo residuos organicos}: frutas, vegetales, cafe, e incluso estiercol.
      \end{itemize}
    \end{column}
    \begin{column}{0.40\textwidth}
      \centering
      \begin{tikzpicture}[scale=0.8]
        \draw[fill=marronlarva,rounded corners] (0,0) ellipse (1.2 and 0.4);
        \draw[fill=marronlarva!70] (-0.8,0.2) -- (0.8,0.2) -- (0.6,1.5) -- (-0.6,1.5) -- cycle;
        \draw[fill=marronlarva!50] (-0.4,1.5) -- (0.4,1.5) -- (0.2,2.2) -- (-0.2,2.2) -- cycle;
        \node at (0,-0.6) {\footnotesize\textit{Larva de BSF}};
      \end{tikzpicture}
    \end{column}
  \end{columns}
\end{frame}

\begin{frame}{Un ciclo de vida veloz}
  \centering
  \begin{tikzpicture}[scale=0.9]
    \node[draw,circle,fill=white,minimum size=1cm] (huevo) at (0,0) {\footnotesize Huevo};
    \node[below] at (huevo.south) {\tiny 4 dias};
    \node[draw,rounded corners,fill=marronlarva!30,minimum width=1.5cm,minimum height=0.8cm] (larva) at (3,0) {\footnotesize Larva};
    \node[below] at (larva.south) {\tiny 14 dias};
    \node[draw,diamond,fill=gray!30,minimum size=1cm] (prepupa) at (6,0) {\footnotesize Prepupa};
    \node[below] at (prepupa.south) {\tiny migra};
    \node[draw,regular polygon,regular polygon sides=6,fill=verdehoja!20,minimum size=1cm] (pupa) at (9,0) {\footnotesize Pupa};
    \node[below] at (pupa.south) {\tiny 14 dias};
    \node[draw,circle,fill=azulcielo!30,minimum size=1cm] (adulto) at (12,0) {\footnotesize Adulto};
    \node[below] at (adulto.south) {\tiny 5--8 dias};
    \draw[->,thick,verdehoja] (huevo) -- (larva);
    \draw[->,thick,verdehoja] (larva) -- (prepupa);
    \draw[->,thick,verdehoja] (prepupa) -- (pupa);
    \draw[->,thick,verdehoja] (pupa) -- (adulto);
    \draw[->,thick,verdehoja,dashed] (adulto.north) to[out=120,in=60] (huevo.north);
  \end{tikzpicture}
  \vspace{0.5cm}
  \begin{block}{El dato clave}
    Solo la \textbf{larva} come. El adulto vive solo para reproducirse.\\
    \textbf{El insecto productivo = la larva.}
  \end{block}
\end{frame}

\begin{frame}{Por que nos interesa?}
  \begin{columns}
    \begin{column}{0.48\textwidth}
      \textbf{El problema}
      \begin{itemize}
        \item Colombia genera $\sim$12 millones de toneladas de residuos organicos al ano.
        \item Gran parte va a rellenos sanitarios, generando metano (CH$_4$).
        \item Necesitamos alternativas que \textbf{valoricen} en vez de \textbf{contaminar}.
      \end{itemize}
    \end{column}
    \begin{column}{0.48\textwidth}
      \textbf{La solucion}
      \begin{itemize}
        \item Las larvas convierten residuos en \textbf{biomasa rica en proteina} (42\% PB).
        \item El residuo (frass) es un \textbf{excelente fertilizante organico}.
        \item Reduccion de emisiones de GEI vs. descomposicion natural.
      \end{itemize}
    \end{column}
  \end{columns}
  \vspace{0.4cm}
  \centering
  \begin{tikzpicture}
    \node[draw,rounded corners,fill=verdehoja!20] (residuo) at (0,0) {Residuos};
    \node[draw,rounded corners,fill=marronlarva!30] (larva) at (4,0) {Larvas BSF};
    \node[draw,rounded corners,fill=azulcielo!20] (proteina) at (8,0.8) {Proteina animal};
    \node[draw,rounded corners,fill=verdehoja!30] (frass) at (8,-0.8) {Fertilizante};
    \draw[->,thick] (residuo) -- (larva) node[midway,above] {\footnotesize bioconversion};
    \draw[->,thick] (larva) -- (proteina) node[midway,above,sloped] {\footnotesize harina};
    \draw[->,thick] (larva) -- (frass) node[midway,below,sloped] {\footnotesize frass};
  \end{tikzpicture}
\end{frame}

\begin{frame}{El reto: de donde sale el CO$_2$?}
  \begin{columns}
    \begin{column}{0.55\textwidth}
      En la camara de cria medimos CO$_2$, pero este viene de \textbf{dos fuentes}:
      \vspace{0.3cm}
      \begin{enumerate}
        \item \textbf{La larva:} respira para crecer, moverse y mantenerse viva.
        \item \textbf{El lecho:} bacterias y hongos degradan lo que la larva no comio.
      \end{enumerate}
      \vspace{0.3cm}
      El sensor \textbf{no las separa}.\\
      Como sabemos cuanto le toca a cada una?
    \end{column}
    \begin{column}{0.40\textwidth}
      \centering
      \begin{tikzpicture}[scale=0.7]
        \draw[thick,rounded corners,fill=gray!10] (-2,-1.5) rectangle (2,2);
        \node at (0,1.6) {\footnotesize Camara de cria};
        \draw[fill=marronlarva!50] (-0.8,0) ellipse (0.5 and 0.2);
        \node at (-0.8,-0.5) {\tiny larva};
        \draw[fill=verdehoja!30] (0.8,0) circle (0.4);
        \node at (0.8,-0.6) {\tiny bacterias};
        \draw[->,thick,red] (0,2.2) -- (0,1.5);
        \node at (0,2.5) {\tiny CO$_2$};
        \draw[->,thick,blue] (-0.8,0.3) to[out=60,in=120] (0,1.5);
        \draw[->,thick,green!70!black] (0.8,0.4) to[out=60,in=60] (0,1.5);
      \end{tikzpicture}
    \end{column}
  \end{columns}
\end{frame}

\begin{frame}{Que es un modelo DEB?}
  \begin{block}{Dynamic Energy Budget}
    Una ``hoja de balance'' matematica que rastrea donde va la energia del alimento:
  \end{block}
  \vspace{0.3cm}
  \centering
  \begin{tikzpicture}[scale=0.85]
    \node[draw,rounded corners,fill=yellow!30] (alimento) at (0,0) {Alimento};
    \node[draw,rounded corners,fill=orange!30] (reservas) at (3.5,0) {Reservas (lipidos)};
    \node[draw,rounded corners,fill=verdehoja!30] (estructura) at (7,0.8) {Estructura (crecer)};
    \node[draw,rounded corners,fill=red!20] (mantenimiento) at (7,-0.8) {Mantenimiento (respirar)};
    \node[draw,rounded corners,fill=gray!20] (lipidos) at (10.5,0) {Lipidos (ahorro)};
    \draw[->,thick] (alimento) -- (reservas) node[midway,above] {\footnotesize asimila};
    \draw[->,thick] (reservas) -- (estructura) node[midway,above,sloped] {\footnotesize $\kappa$};
    \draw[->,thick] (reservas) -- (mantenimiento) node[midway,below,sloped] {\footnotesize $1-\kappa$};
    \draw[->,thick] (reservas) -- (lipidos) node[midway,above,sloped] {\footnotesize excedente};
  \end{tikzpicture}
  \vspace{0.3cm}
  \begin{itemize}
    \item Desarrollado por Bas Kooijman (Vrije Universiteit Amsterdam).
    \item Aplicado a peces, crustaceos, ahora a insectos (Eriksen 2022).
    \item Nosotros lo extendimos para incluir el \textbf{lecho bacteriano}.
  \end{itemize}
\end{frame}

\begin{frame}{Nuestro aporte: un gemelo digital del batch}
  \begin{columns}
    \begin{column}{0.50\textwidth}
      \textbf{Que hicimos:}
      \begin{itemize}
        \item Disenamos un batch controlado con 2 dietas (alta y baja calidad).
        \item Medimos CO$_2$ cada hora, biomasa cada dia, carbono del sustrato.
        \item Calibramos un modelo que separa larva $\leftrightarrow$ lecho.
      \end{itemize}
      \vspace{0.3cm}
      \textbf{Que obtuvimos:}
      \begin{itemize}
        \item El nucleo Eriksen reproduce el crecimiento larvario.
        \item El lecho explica el CO$_2$ cuando \textbf{no hay larvas} (control).
        \item El pastoreo larvario es indispensable: sin el, el lecho se desborda.
      \end{itemize}
    \end{column}
    \begin{column}{0.45\textwidth}
      \centering
      \begin{tikzpicture}[scale=0.7]
        \draw[thick,rounded corners,fill=gray!10] (-2.5,-2) rectangle (2.5,2.5);
        \node at (0,2) {\footnotesize \textbf{Gemelo digital}};
        \draw[fill=marronlarva!40] (-1,0.5) ellipse (0.6 and 0.3);
        \node at (-1,0) {\tiny Larva (DEB)};
        \draw[fill=verdehoja!30] (1,0.5) circle (0.4);
        \node at (1,0) {\tiny Lecho (Monod)};
        \draw[<->,thick,marronlarva] (-0.3,0.5) -- (0.5,0.5);
        \node at (0.1,0.8) {\tiny pastoreo};
        \draw[->,thick,blue] (0,-1.5) -- (0,-0.3);
        \node at (0.8,-1.2) {\tiny CO$_2$ total};
        \node[draw,rounded corners,fill=white] at (0,-1.8) {\tiny Prediccion};
      \end{tikzpicture}
    \end{column}
  \end{columns}
\end{frame}

\begin{frame}{Resultados: funciona?}
  \centering
  \textbf{Validacion externa: nuestro modelo predice datos de otros laboratorios?}
  \vspace{0.4cm}
  \begin{columns}
    \begin{column}{0.48\textwidth}
      \begin{block}{Ermolaev (Suecia, 2019)}
        700 larvas, residuos de cocina, 28$^\circ$C.
        \vspace{0.2cm}
        
        $\Rightarrow$ Error en acumulado: \textbf{menos del 3\%}.
        \vspace{0.1cm}
        
        El nucleo DEB \textbf{generaliza} sin reajuste.
      \end{block}
    \end{column}
    \begin{column}{0.48\textwidth}
      \begin{block}{Schon (Alemania, 2026)}
        Dietas de distinta calidad proteica.
        \vspace{0.2cm}
        
        $\Rightarrow$ Ordenamiento correcto en 4/4 dietas.
        \vspace{0.1cm}
        
        Limite: en dieta muy pobre, el nucleo \textbf{se satura}.
      \end{block}
    \end{column}
  \end{columns}
  \vspace{0.4cm}
  \textcolor{verdehoja}{\textbf{Narrativa honesta:}} el modelo funciona bien en condiciones normales, pero declaramos sus limites.
\end{frame}

\begin{frame}{Y esto para que sirve en el Valle del Cauca?}
  \begin{columns}
    \begin{column}{0.55\textwidth}
      \textbf{Residuos disponibles:}
      \begin{itemize}
        \item \textbf{Cafe:} mucilago, cascarilla (miles de toneladas/ano).
        \item \textbf{Cana:} bagazo, cachaza (ya se usa, pero sobra).
        \item \textbf{Frutas:} banano, maracuya, aguacate (rechazos de exportacion).
        \item \textbf{Mercados:} residuos organicos de Plaza de Mercado (Cali, Palmira).
      \end{itemize}
      \vspace{0.2cm}
      \textbf{Productos:}
      \begin{itemize}
        \item Harina de larva $\rightarrow$ alimento para peces y aves.
        \item Frass $\rightarrow$ fertilizante organico certificado.
      \end{itemize}
    \end{column}
    \begin{column}{0.40\textwidth}
      \centering
      \begin{tikzpicture}[scale=0.65]
        \node[draw,rounded corners,fill=brown!20] (cafe) at (0,2) {Cafe};
        \node[draw,rounded corners,fill=green!20] (cana) at (0,0.8) {Cana};
        \node[draw,rounded corners,fill=yellow!20] (fruta) at (0,-0.4) {Frutas};
        \node[draw,rounded corners,fill=gray!20] (mercado) at (0,-1.6) {Mercados};
        \node[draw,rounded corners,fill=marronlarva!40,minimum width=1.5cm] (bsf) at (4,0.2) {BSF};
        \node[draw,rounded corners,fill=red!20] (proteina) at (8,1) {Proteina};
        \node[draw,rounded corners,fill=verdehoja!30] (frass) at (8,-0.6) {Fertilizante};
        \draw[->,thick] (cafe) -- (bsf);
        \draw[->,thick] (cana) -- (bsf);
        \draw[->,thick] (fruta) -- (bsf);
        \draw[->,thick] (mercado) -- (bsf);
        \draw[->,thick] (bsf) -- (proteina);
        \draw[->,thick] (bsf) -- (frass);
      \end{tikzpicture}
    \end{column}
  \end{columns}
\end{frame}

\begin{frame}{Hacia donde vamos: el gemelo digital IoT}
  \begin{columns}
    \begin{column}{0.55\textwidth}
      \textbf{Vision:} una granja de BSF con sensores en tiempo real.
      \vspace{0.3cm}
      
      \textbf{Que mide el sensor:}
      \begin{itemize}
        \item CO$_2$, temperatura, humedad del lecho.
        \item Peso de la camara (biomasa estimada).
      \end{itemize}
      \vspace{0.2cm}
      
      \textbf{Que predice el modelo:}
      \begin{itemize}
        \item La larva esta creciendo bien? Necesita mas alimento?
        \item El lecho se esta desbordando? Hay que ajustar humedad?
        \item Cuando cosechar? (antes de que empiece la migracion).
      \end{itemize}
    \end{column}
    \begin{column}{0.40\textwidth}
      \centering
      \begin{tikzpicture}[scale=0.7]
        \draw[thick,rounded corners,fill=blue!5] (-2,-2) rectangle (2,2.5);
        \node at (0,2) {\footnotesize \textbf{Dashboard}};
        \draw[fill=green!40] (-1.5,0.5) rectangle (-0.5,1.2);
        \node at (-1,0.85) {\tiny CO$_2$};
        \draw[fill=orange!40] (-0.3,0.5) rectangle (0.7,1.2);
        \node at (0.2,0.85) {\tiny T$^\circ$};
        \draw[fill=purple!30] (0.9,0.5) rectangle (1.9,1.2);
        \node at (1.4,0.85) {\tiny H\%};
        \draw[fill=gray!30] (-1.5,-0.5) rectangle (1.9,0.2);
        \node at (0.2,-0.15) {\tiny Prediccion de biomasa};
        \draw[fill=red!20] (-1.5,-1.3) rectangle (1.9,-0.6);
        \node at (0.2,-0.95) {\tiny Alerta: ajustar alimento};
      \end{tikzpicture}
    \end{column}
  \end{columns}
\end{frame}

\begin{frame}{Conclusiones}
  \begin{enumerate}
    \item \textit{Hermetia illucens} es una \textbf{alternativa real} para valorizar residuos organicos del Valle del Cauca.
    \item El \textbf{modelo DEB} permite rastrear la energia del alimento $\rightarrow$ crecimiento $\rightarrow$ emision, con cierre de carbono.
    \item Nuestro modelo \textbf{hibrido} (larva DEB + lecho Monod) separa las fuentes de CO$_2$ por primera vez en BSF.
    \item El \textbf{pastoreo bacteriano} es clave: las larvas no solo comen el sustrato, tambien comen las bacterias que lo degradan.
    \item El \textbf{gemelo digital} esta al alcance: simulacion en $<1$ segundo, calibrable con datos de los primeros dias.
  \end{enumerate}
\end{frame}

\begin{frame}
  \vspace{1.5cm}
  \centering
  {\Huge\textcolor{verdehoja}{Gracias}}\\[0.8cm]
  
  {\normalsize\textcolor{marronlarva}{Preguntas?}}\\[1cm]
  
  {\small \textit{jjlealg@unal.edu.co}}\\[0.2cm]
  {\footnotesize\texttt{github.com/asoleal/tesis-bsf-co2}}\\[0.5cm]
  
  \begin{tikzpicture}
    \draw[fill=verdehoja,rounded corners] (-0.5,-0.2) rectangle (0.5,0.2);
    \node[white] at (0,0) {\tiny BSF};
  \end{tikzpicture}
\end{frame}

\end{document}
